\section{Discussion}
\label{sec:discussion}

\subsection{Predictive Power of Cross-Sectional Analysis}
\label{sec:disc_predictive}

The most striking finding is the perfect concordance between cross-sectional predictions and causal validation: all six predictions derived from observational behavioral profiling are confirmed by the natural experiment (Table~\ref{tab:validation}). This concordance has methodological significance beyond the specific policy context. It demonstrates that cross-sectional multi-outcome profiling---comparing behavioral outcomes across concurrent policy regimes on identical hardware---can accurately forecast which margins will respond to policy changes and which will not.

The key discriminator is outcome sensitivity to mode. Outcomes with large GEE mode effects and high SHAP importance (harsh acceleration $\beta = -0.61$, harsh deceleration $\beta = -0.63$) responded causally to the TUB ban. Outcomes with small mode effects and low SHAP importance (speed CV $\beta = -0.02$, SHAP $= 0.03$) showed null causal effects. This pattern suggests a practical rule: the cross-sectional mode effect size serves as a leading indicator for the causal policy response. Cities or operators considering governor mandates could use observational behavioral profiling as a pre-implementation assessment tool---characterizing which safety, demand, and dynamics margins are mode-sensitive before the policy change occurs.

\subsection{Selective Response Across Margins}
\label{sec:disc_selective}

Speed governance produces a remarkably \textit{selective} behavioral response. The single policy lever (capping maximum speed) causally reduces acceleration aggressiveness---harsh acceleration and deceleration events drop significantly after the TUB ban, surviving Bonferroni correction---while leaving demand (trip counts, active users), riding dynamics (cruising patterns, stop-and-go behavior), and speed variability completely unaffected. The predict-then-validate framework makes this selectivity more than a post hoc observation: the cross-sectional evidence \textit{predicted} exactly which margins would respond, transforming a descriptive finding into a validated forecasting result.

\subsection{No Compensation: An Anti-Peltzman Result}
\label{sec:disc_peltzman}

The Peltzman hypothesis \citep{peltzman1975effects} and risk homeostasis theory \citep{wilde1982theory} predict that safety interventions may be partially offset by behavioral compensation on unconstrained margins. Our results provide the first multi-margin test of this hypothesis for e-scooter speed governance, and the answer is unambiguous: \textit{zero compensation on any measured margin}.

The mode-switcher design cleanly separates composition from compensation---a distinction that prior aggregate analyses cannot make. The STD/ECO-only DiD shows significant \textit{increases} in harsh events and speed after the TUB ban---an apparent Peltzman effect. But the within-user analysis reveals all Cohen's $d$ values below 0.14 (negligible). Former TUB users entering the STD/ECO pool are inherently more aggressive riders (pre-existing speed gap of 0.60~km/h), producing a compositional shift in the STD/ECO pool, not behavioral adaptation. Without the mode-switcher decomposition, a researcher observing only aggregate STD/ECO outcomes would erroneously conclude that the TUB ban \textit{induced} compensation---precisely the type of ecological fallacy that the Peltzman debate has been unable to resolve for decades.

This finding aligns with \citet{hedlund2000risky}, who argued that behavioral compensation is most plausible when the intervention is \textit{salient} (the user knows protection has increased) and the regulated behavior is \textit{voluntary}. Speed governors fail both conditions: they operate invisibly (riders on governed modes may not perceive any safety ``benefit'' from the cap), and the speed constraint is involuntary (the vehicle physically cannot exceed the limit). The seatbelt analogy is instructive: \citet{cohen2003effects} found that mandatory seatbelt laws reduced occupant fatalities without inducing measurable increases in risky driving, contra Peltzman's original prediction. Our result extends this pattern to a fundamentally different safety technology in a different mode: software-enforced speed limits on shared micromobility.

The null compensation finding has important policy implications: governor mandates need not be tempered by concerns about offsetting behavioral responses. The safety benefit of reduced harsh events is a net gain, not partially offset by compensation on other margins. This is particularly noteworthy because the e-scooter context offered many unconstrained margins on which compensation \textit{could} have manifested: acceleration aggressiveness within governed modes, more erratic speed profiles, longer or more dangerous routes, or reduced trip frequency. None materialized.

\subsection{TUB Adoption as an Escalation Pathway}
\label{sec:disc_escalation}

The experience-mediation analysis reveals a concerning pathway: as riders gain experience, they progressively adopt the ungoverned TUB mode (23\% at trip~1 to 70\% at trip~101+), and this mode adoption mediates 57--70\% of the experience-related increase in harsh events. Within governed modes (STD/ECO), experienced riders show only modest changes in riding behavior. The escalation is primarily a \textit{mode selection} story: experienced riders learn to use the faster option, not that they inherently ride more aggressively.

This mode-selection mechanism has a direct analogue in the graduated driver licensing (GDL) literature. GDL programs reduce novice drivers' crash risk by 20--40\% through staged exposure to high-risk conditions \citep{shope2007graduated}---nighttime driving, passenger restrictions, and highway access are unlocked sequentially as experience accumulates. The parallel to e-scooter speed governance is exact: restricting novice riders to governed modes and allowing access to higher-speed modes only after demonstrated safe riding could interrupt the escalation pathway without imposing permanent restrictions. The survival analysis quantifies the potential impact: non-TUB users have a median time-to-first-speeding of infinity (most never speed), while TUB users reach it at trip~7. A policy that delays TUB access until, say, trip~50 would prevent the early-career speeding events that currently characterize the TUB adoption trajectory.

Importantly, the TUB adoption curve is concave: it rises steeply in the first 20--30 trips before plateauing around 70\%. This suggests a critical window during which mode restriction would be most effective. Riders who are ``nudged'' toward governed modes during this initial period may develop riding habits that persist even after restrictions are lifted---a hypothesis supported by the GDL finding that crash reductions persist beyond the restriction period.

\subsection{Policy Implications}
\label{sec:disc_policy}

Three policy recommendations emerge:

First, \textbf{mandatory speed governors} on all shared e-scooters. The causal evidence shows that governor imposition reduces harsh events with no demand cost (trip counts and active users are unaffected) and no behavioral compensation. The policy is Pareto-improving on measured margins.

Second, \textbf{graduated mode access}. Restricting new riders to governed modes could interrupt the TUB adoption escalation pathway without requiring permanent mode restrictions. As riders demonstrate safe behavior, higher-speed modes could be unlocked---analogous to graduated driver licensing.

Third, \textbf{cross-sectional profiling as pre-implementation assessment}. Before mandating speed governors, regulators can deploy multi-outcome behavioral profiling to forecast which margins will respond and whether demand costs are likely. Our predict-then-validate results show that this approach works: the cross-sectional mode effect size accurately forecasts the causal policy effect across all six tested outcomes. This provides a practical template for evaluating other speed management interventions (geofencing, dynamic speed limits, incentive-based approaches) before implementation.

\subsection{Limitations}
\label{sec:disc_limitations}

Several limitations should be acknowledged regarding data scope, measurement, and econometric identification.

\textit{Data and design.} Our data come from a single operator (Swing) in a single country (South Korea), and generalizability to other operators, vehicle types, and cultural contexts requires further study. Mode self-selection in cross-sectional analyses is addressed by individual fixed effects (absorbing time-invariant user traits) and DiD (exploiting an exogenous policy shock), though time-varying unobservables (e.g., within-day mood or fatigue) remain uncontrolled. The December 2023 analysis includes only one post-treatment month; longer-term effects, including potential delayed compensation, cannot be assessed.

\textit{Measurement.} We analyze behavioral proxies (harsh events, speed variability) rather than crash outcomes; the causal link between reduced harsh events and reduced injuries, while theoretically grounded in the power model \citep{elvik2013reexamination}, is not directly tested. The $\sim$10-second GPS sampling interval smooths rapid acceleration events; higher-frequency kinematic data (e.g., accelerometer) would capture a wider range of aggressive riding behaviors.

\textit{Econometric identification.} The Cox proportional hazards models in the escalation analysis show violations of the PH assumption for tub\_frac in all three endpoints (Schoenfeld residual test $p < 0.001$), meaning that the effect of TUB usage on event hazard is not constant over the trip sequence. Schoenfeld residual--time correlations are uniformly negative for tub\_frac ($\rho = -0.05$ to $-0.15$, all $p < 0.001$), indicating that the TUB hazard ratio \textit{attenuates} over the trip sequence---the strongest TUB effect occurs in early trips and diminishes as riders accumulate experience. This attenuation is substantively plausible: early TUB adoption may represent the highest-risk transition, while experienced TUB users develop greater familiarity with the vehicle's capabilities. The reported HRs should therefore be interpreted as average effects over the trip sequence rather than constant multipliers; the qualitative conclusions (TUB strongly accelerates time-to-first-speeding and time-to-first-harsh-event) are robust to this concern, as confirmed by the non-parametric KM curves. Recent econometrics literature has identified potential bias in TWFE estimators under staggered treatment timing \citep{goodman2021difference}; our design has a single system-wide treatment date, avoiding this concern, but heterogeneity-robust estimators could be applied in future extensions with staggered rollouts. Finally, the parallel trends assumption is supported by the placebo test but the event study shows some pre-trend variation during the February--March ramp-up period; we note that this pre-trend is positive (conservative), meaning our treatment effect estimate may be attenuated.
